% GENERATED FILE. Edit canonical lessons, then run npm run generate:books.
% canonical-source-hash: fnv1a64:ad0b9d73541547f5
% canonical-lessons: FA-C456-peyk, FA-C456-sanduq-e-post, FA-C456-simkart, FA-C456-bongah, FA-C456-tamas-gereftan

\chapter{Couriers, phones and agents}
\label{ch:fa-couriers-phones-and-agents}
\hlchaptermodality{456}

\begin{chapteropening}
\textbf{By the end of this chapter:} \emph{I can use services in town: a courier, a postbox, a SIM card, an estate agency and getting in touch.}

Complete the last lesson of chapter 456: I can use services in town: a courier, a postbox, a SIM card, an estate agency and getting in touch.
\end{chapteropening}

\section[\texorpdfstring{\fa{پیک} (peyk)}{پیک (peyk)}]{\fa{پیک} (peyk) — a courier, a delivery rider}

\label{lesson:FA-C456-peyk}



\begin{quote}
Before the new one: say the Persian for a cultural centre, then the Persian for to close.
\end{quote}

\subsection*{You'll want to know: \fa{پیک}}
\textbf{\fa{پیک}} (\emph{peyk}) — \textquotedblleft{}a courier, a delivery rider\textquotedblright{}.

\begin{quote}
\fa{پیک} \fa{غذا} \fa{را} \fa{جلوی} \fa{در} \fa{آورد}.
\end{quote}

\begin{quote}
\emph{Peyk ghazâ râ jolow-ye dar âvard.}
\end{quote}

\begin{quote}
The delivery rider brought the food to the door.
\end{quote}

\emph{Peyk-e motori} — a motorbike courier.

\begin{cousinweb}
\emph{peyk} is inherited: EtymDB traces it to Middle Persian \emph{payg}, which is usually glossed 'a messenger, a runner'.
\end{cousinweb}

\begin{grammarlens}[title={What you've built}]
One more word for telling, explaining and describing at B1.
\end{grammarlens}

\subsection*{Your turn}
\begin{itemize}
\raggedright
  \item \emph{Say it:} \emph{peyk}
  \item \emph{Say it:} \emph{peyk}, once more
  \item \emph{Say it:} a sentence of your own with \emph{peyk}
  \item \emph{Recall:} say the Persian for a cultural centre, then the Persian for to close, then say \emph{peyk} again
\end{itemize}

\begin{tcolorbox}[breakable,colback=teal!4,colframe=teal!35!black,title={Before you move on}]
What does \fa{پیک} mean? (\textquotedblleft{}A courier, a delivery rider\textquotedblright{}.) Say it once more.
\end{tcolorbox}
\section[\texorpdfstring{\fa{صندوق} \fa{پست} (sanduq-e post)}{صندوق پست (sanduq-e post)}]{\fa{صندوق} \fa{پست} (sanduq-e post) — a postbox (US: a mailbox)}

\label{lesson:FA-C456-sanduq-e-post}



\begin{quote}
Before the new one: say the Persian for to close, then the Persian for a courier.
\end{quote}

\subsection*{You'll want to know: \fa{صندوق} \fa{پست}}
\textbf{\fa{صندوق} \fa{پست}} (\emph{sanduq-e post}) — \textquotedblleft{}a postbox (US: a mailbox)\textquotedblright{}.

\begin{quote}
\fa{نامه} \fa{را} \fa{در} \fa{صندوق} \fa{پست} \fa{انداختم}.
\end{quote}

\begin{quote}
\emph{Nâme râ dar sanduq-e post andâkhtam.}
\end{quote}

\begin{quote}
I posted the letter in the postbox.
\end{quote}

\emph{sanduq} 'box' + \emph{post} 'post'.

\begin{cousinweb}
\emph{sanduq-e post} is built from \emph{sanduq} 'chest, box', borrowed from Arabic \emph{ṣundūq} (EtymDB), and \emph{post}, a loan usually traced to French \emph{poste}.
\end{cousinweb}

\begin{grammarlens}[title={What you've built}]
One more word for telling, explaining and describing at B1.
\end{grammarlens}

\subsection*{Your turn}
\begin{itemize}
\raggedright
  \item \emph{Say it:} \emph{sanduq-e post}
  \item \emph{Say it:} \emph{sanduq-e post}, once more
  \item \emph{Say it:} a sentence of your own with \emph{sanduq-e post}
  \item \emph{Recall:} say the Persian for to close, then the Persian for a courier, then say \emph{sanduq-e post} again
\end{itemize}

\begin{tcolorbox}[breakable,colback=teal!4,colframe=teal!35!black,title={Before you move on}]
What does \fa{صندوق} \fa{پست} mean? (\textquotedblleft{}A postbox (US: a mailbox)\textquotedblright{}.) Say it once more.
\end{tcolorbox}
\section[\texorpdfstring{\fa{سیم} \fa{کارت} (simkârt)}{سیم کارت (simkârt)}]{\fa{سیم} \fa{کارت} (simkârt) — a SIM card}

\label{lesson:FA-C456-simkart}



\begin{quote}
Before the new one: say the Persian for a courier, then the Persian for a postbox.
\end{quote}

\subsection*{You'll want to know: \fa{سیم} \fa{کارت}}
\textbf{\fa{سیم} \fa{کارت}} (\emph{simkârt}) — \textquotedblleft{}a SIM card\textquotedblright{}.

\begin{quote}
\fa{در} \fa{فرودگاه} \fa{یک} \fa{سیم} \fa{کارت} \fa{تازه} \fa{خریدم}.
\end{quote}

\begin{quote}
\emph{Dar forudgâh yek simkârt-e tâze kharidam.}
\end{quote}

\begin{quote}
At the airport I bought a new SIM card.
\end{quote}

Written like native \emph{sim} 'wire' + \emph{kârt}, but taken from English.

\begin{cousinweb}
\emph{simkârt} was borrowed from English \emph{SIM card}. It is spelled like native \emph{sim} 'wire', which EtymDB traces to Middle Persian, and \emph{kârt} 'card', from French \emph{carte}.
\end{cousinweb}

\begin{grammarlens}[title={What you've built}]
One more word for telling, explaining and describing at B1.
\end{grammarlens}

\subsection*{Your turn}
\begin{itemize}
\raggedright
  \item \emph{Say it:} \emph{simkârt}
  \item \emph{Say it:} \emph{simkârt}, once more
  \item \emph{Say it:} a sentence of your own with \emph{simkârt}
  \item \emph{Recall:} say the Persian for a courier, then the Persian for a postbox, then say \emph{simkârt} again
\end{itemize}

\begin{tcolorbox}[breakable,colback=teal!4,colframe=teal!35!black,title={Before you move on}]
What does \fa{سیم} \fa{کارت} mean? (\textquotedblleft{}A SIM card\textquotedblright{}.) Say it once more.
\end{tcolorbox}
\section[\texorpdfstring{\fa{بنگاه} (bongâh)}{بنگاه (bongâh)}]{\fa{بنگاه} (bongâh) — an estate agency; a firm, an agency}

\label{lesson:FA-C456-bongah}



\begin{quote}
Before the new one: say the Persian for a postbox, then the Persian for a SIM card.
\end{quote}

\subsection*{You'll want to know: \fa{بنگاه}}
\textbf{\fa{بنگاه}} (\emph{bongâh}) — \textquotedblleft{}an estate agency; a firm, an agency\textquotedblright{}.

\begin{quote}
\fa{از} \fa{طریق} \fa{بنگاه} \fa{یک} \fa{آپارتمان} \fa{پیدا} \fa{کردیم}.
\end{quote}

\begin{quote}
\emph{Az tariq-e bongâh yek âpârtemân peydâ kardim.}
\end{quote}

\begin{quote}
We found a flat through an estate agency.
\end{quote}

In full \emph{bongâh-e moshâver-e amlâk}. Also a firm: \emph{bongâh-e bârbari} — a removals firm.

\begin{cousinweb}
\emph{bongâh} is usually analysed as native \emph{bon} 'base, foundation' and \emph{gâh} 'place'. EtymDB does not list it.
\end{cousinweb}

\begin{grammarlens}[title={What you've built}]
One more word for telling, explaining and describing at B1.
\end{grammarlens}

\subsection*{Your turn}
\begin{itemize}
\raggedright
  \item \emph{Say it:} \emph{bongâh}
  \item \emph{Say it:} \emph{bongâh}, once more
  \item \emph{Say it:} a sentence of your own with \emph{bongâh}
  \item \emph{Recall:} say the Persian for a postbox, then the Persian for a SIM card, then say \emph{bongâh} again
\end{itemize}

\begin{tcolorbox}[breakable,colback=teal!4,colframe=teal!35!black,title={Before you move on}]
What does \fa{بنگاه} mean? (\textquotedblleft{}An estate agency; a firm, an agency\textquotedblright{}.) Say it once more.
\end{tcolorbox}
\section[\texorpdfstring{\fa{تماس} \fa{گرفتن} (tamâs gereftan)}{تماس گرفتن (tamâs gereftan)}]{\fa{تماس} \fa{گرفتن} (tamâs gereftan) — to get in touch, to call (someone)}

\label{lesson:FA-C456-tamas-gereftan}



\begin{quote}
Before the new one: say the Persian for a SIM card, then the Persian for an estate agency.
\end{quote}

\subsection*{You'll want to know: \fa{تماس} \fa{گرفتن}}
\textbf{\fa{تماس} \fa{گرفتن}} (\emph{tamâs gereftan}) — \textquotedblleft{}to get in touch, to call (someone)\textquotedblright{}.

\begin{quote}
\fa{با} \fa{هتل} \fa{تماس} \fa{گرفتیم} \fa{و} \fa{اتاق} \fa{را} \fa{رزرو} \fa{کردیم}.
\end{quote}

\begin{quote}
\emph{Bâ hotel tamâs gereftim va otâq râ rezerv kardim.}
\end{quote}

\begin{quote}
We called the hotel and booked the room.
\end{quote}

The person called takes \emph{bâ}. More formal than \emph{zang zadan}.

\begin{cousinweb}
\emph{tamâs gereftan} is built from \emph{tamâs} 'contact', a loan usually traced to Arabic \emph{tamāss} 'touching one another', and the native verb \emph{gereftan} 'to take'. EtymDB does not list \emph{tamâs}.
\end{cousinweb}

\begin{grammarlens}[title={What you've built}]
One more word for telling, explaining and describing at B1.
\end{grammarlens}

\subsection*{Your turn}
\begin{itemize}
\raggedright
  \item \emph{Say it:} \emph{tamâs gereftan}
  \item \emph{Say it:} \emph{tamâs gereftan}, once more
  \item \emph{Say it:} a sentence of your own with \emph{tamâs gereftan}
  \item \emph{Recall:} say the Persian for a SIM card, then the Persian for an estate agency, then say \emph{tamâs gereftan} again
\end{itemize}

\begin{tcolorbox}[breakable,colback=teal!4,colframe=teal!35!black,title={Before you move on}]
What does \fa{تماس} \fa{گرفتن} mean? (\textquotedblleft{}To get in touch, to call (someone)\textquotedblright{}.) Say it once more.
\end{tcolorbox}
